% Shared result, cited from solutions with \appendixref{darboux-theorem}.
\begin{appendixitem}{Darboux's Theorem}
    Let
    \[
        \mathcal{I} := \bigcup_{a, b \in \mathbb{R}} \left\{ \left[a;b\right], \left]a;b\right[,
        \left[a;b\right[, \left]a;b\right], \left[a;+\infty\right[, \left]a;+\infty\right[,
        \left]-\infty;b\right], \left]-\infty;b\right[ \right\} \cup \left\{ \mathbb{R} \right\}
    \]
    be the set of intervals of \( \mathbb{R} \), let \( I \in \mathcal{I} \), and let \( \varphi
    \colon I \to \mathbb{R} \) be a map. Recall that a subset \( S \) of \( \mathbb{R} \) is an
    interval, i.e., \( S \in \mathcal{I} \), if and only if it contains every real number lying
    between two of its elements, that is, \( \left[a;b\right] \subset S \) for all \( a, b \in S \),
    and also if and only if it is connected for the usual topology. Say that \( \varphi \) has the
    Darboux property, or intermediate value property, if it sends connected sets (of \( I \)) to
    connected sets, i.e., \( \varphi\left[J\right] \in \mathcal{I} \) for every \( J \in
    \mathcal{I} \) with \( J \subset I \). Recall also that \( \varphi \) is differentiable at
    \( x \in I \) if \( \frac{\varphi\left(t\right)-\varphi\left(x\right)}{t-x} \) has a limit
    \( \ell \in \mathbb{R} \) as \( t \to x \) with \( t \in I \setminus
    \left\{x\right\} \)\footnote{That is, for every open set \( U \) of \( \mathbb{R} \) containing
    \( \ell \), there exists an open set \( E \) of \( I \) containing \( x \) such that
    \( \frac{\varphi\left(t\right)-\varphi\left(x\right)}{t-x} \in U \) for all \( t \in E \setminus
    \left\{x\right\} \). When \( I \) has at least two points, \( x \) is a limit point of
    \( I \setminus \left\{x\right\} \), so that this limit is unique (\( \mathbb{R} \) is
    Hausdorff). Equivalently, \( \varphi \) is differentiable at \( x \) with
    \( \varphi'\left(x\right) = \ell \) if and only if
    \( \frac{\varphi\left(t_n\right)-\varphi\left(x\right)}{t_n-x}
    \overset{n \to +\infty}{\longrightarrow} \ell \) for every sequence \( t_{-} \in \left(I
    \setminus \left\{x\right\}\right)^{\mathbb{N}_{\geq 1}} \) with \( t_n
    \overset{n \to +\infty}{\longrightarrow} x \). The direct implication holds for limits in any
    topological space. The converse only uses that \( x \) has a countable basis of neighbourhoods
    \( V_1 \supset V_2 \supset \cdots \) in \( I \) (for instance \( V_n :=
    \left]x-\frac{1}{n};x+\frac{1}{n}\right[ \cap I \)), i.e., that \( I \) is first countable at
    \( x \). Indeed, by contraposition, if the quotient did not tend to \( \ell \) (in the first
    sense), there would exist an open set \( U \) of \( \mathbb{R} \) containing \( \ell \) and, for
    each \( n \in \mathbb{N}_{\geq 1} \), some \( t_n \in V_n \setminus \left\{x\right\} \) with
    \( \frac{\varphi\left(t_n\right)-\varphi\left(x\right)}{t_n-x} \notin U \) (by the axiom of
    countable choice), so that \( t_n \overset{n \to +\infty}{\longrightarrow} x \) while the
    quotients \( \frac{\varphi\left(t_n\right)-\varphi\left(x\right)}{t_n-x} \) do not tend to
    \( \ell \).}, denoted by \( \varphi'\left(x\right) \). In particular, at an endpoint of \( I \)
    that lies in \( I \), this is a one-sided derivative.
    \begin{theorem}[Darboux]
        Let \( I \in \mathcal{I} \) and \( f \colon I \to \mathbb{R} \) be a differentiable map.
        Then \( f' \) has the Darboux property, i.e., \( f'\left[J\right] \) is an interval for
        every interval \( J \subset I \).
    \end{theorem}
    \begin{proof}
        Fix \( J \in \mathcal{I} \) with \( J \subset I \). We show that \( f'\left[J\right] \) is
        an interval. By the characterisation of the intervals recalled above, this is equivalent to
        showing that \( \left[u;v\right] \subset f'\left[J\right] \) for all \( u, v \in
        f'\left[J\right] \). If \( u > v \), then \( \left[u;v\right] = \emptyset \), and if
        \( u = v \), then \( \left[u;v\right] = \left\{u\right\} \subset f'\left[J\right] \), so
        that it is sufficient to treat the case \( u < v \). Clearly, it is sufficient to show that
        \( \left[f'\left(a\right);f'\left(b\right)\right] \subset f'\left[J\right] \) for all
        \( a, b \in J \) with \( f'\left(a\right) < f'\left(b\right) \) (since any such \( u, v \)
        have preimages, and any such preimages will do). As \( f'\left(a\right) \) and
        \( f'\left(b\right) \) themselves lie in \( f'\left[J\right] \), we need to show that
        \( \left]f'\left(a\right);f'\left(b\right)\right[ \subset f'\left[J\right] \). It is then
        sufficient to show the stronger statement that, for all \( y \in
        \left]f'\left(a\right);f'\left(b\right)\right[ \), there exists \( c \) strictly between
        \( a \) and \( b \) (i.e., \( c \in \left]a;b\right[ \cup \left]b;a\right[ \)) such that
        \( f'\left(c\right) = y \), since such a \( c \) lies in \( J \) (\( J \) is an interval
        containing \( a \) and \( b \)).
        \\\\
        Towards this goal, fix \( a, b \in J \) with \( f'\left(a\right) < f'\left(b\right) \) (in
        particular, \( a \neq b \)) and \( y \in \left]f'\left(a\right);f'\left(b\right)\right[ \).
        It is clear that replacing \( \left(f, y, a, b\right) \) by \( \left(-f, -y, b, a\right) \)
        keeps the hypotheses (\( -f \) is differentiable with derivative \( -f' \), and
        \( -f'\left(b\right) < -y < -f'\left(a\right) \)) and the conclusion (a point is in
        \( \left]a;b\right[ \cup \left]b;a\right[ \) exactly when it is in \( \left]b;a\right[ \cup
        \left]a;b\right[ \), and \( -f'\left(c\right) = -y \) if and only if \( f'\left(c\right) =
        y \)). Hence, making this replacement if \( a > b \), we may assume that \( a < b \), so
        that \( f'\left(a\right) < y < f'\left(b\right) \) and we look for \( c \in
        \left]a;b\right[ \) such that \( f'\left(c\right) = y \).
        \\\\
        Define \( g \colon \left[a;b\right] \to \mathbb{R} \), \( x \mapsto f\left(x\right)-yx \),
        which is differentiable (as \( f \) is, on \( \left[a;b\right] \subset I \)) with
        \( g' = f'-y \), so that \( g'\left(a\right) < 0 < g'\left(b\right) \), and it is sufficient
        to show that \( g'\left(c\right) = 0 \) for some \( c \in \left]a;b\right[ \). As \( g \) is
        continuous on the compact interval \( \left[a;b\right] \), it attains its minimum at some
        \( c \in \left[a;b\right] \). Since \( \frac{g\left(x\right)-g\left(a\right)}{x-a}
        \overset{x \to a}{\longrightarrow} g'\left(a\right) < 0 \), there exists \( x \in
        \left]a;b\right[ \) with \( \frac{g\left(x\right)-g\left(a\right)}{x-a} < 0 \), i.e.,
        multiplying by \( x-a > 0 \), \( g\left(x\right) < g\left(a\right) \), so that \( c \neq
        a \) (as \( g\left(c\right) \leq g\left(x\right) < g\left(a\right) \)). Likewise, since
        \( \frac{g\left(x\right)-g\left(b\right)}{x-b} \overset{x \to b}{\longrightarrow}
        g'\left(b\right) > 0 \), there exists \( x \in \left]a;b\right[ \) with
        \( \frac{g\left(x\right)-g\left(b\right)}{x-b} > 0 \), i.e., multiplying by \( x-b < 0 \),
        \( g\left(x\right) < g\left(b\right) \), so that \( c \neq b \) (again, as
        \( g\left(c\right) \leq g\left(x\right) < g\left(b\right) \)). Hence, \( c \in
        \left]a;b\right[ \) is a minimum of \( g \) at an interior point of \( \left[a;b\right] \),
        where \( g \) is differentiable, so that \( g'\left(c\right) = 0 \) by Fermat's theorem, and
        this concludes.
    \end{proof}
\end{appendixitem}
